% page 293 of the facsimile (image p-292 of the scan)
% block 167.6 x 233.3 mm ; left margin 34.4 mm
\pgc{12.700mm}{0.000mm}{
\l{\cel{54}285}
\l{}
\l{}
\l{kolektoro. (\sou{tekn.})\cel{2}(\sou{elektro}) Parto di elektro-mashino, ek asem-}
\l{\cel{3}blitaro de lami ek metalo elektrokonduktiva (kupro, precipue),}
\l{\cel{3}dispozita periferie di cilindro, ed inter-izolita per lami}
\l{\cel{3}de mikao. - DEFIRS.}
\l{}
\l{\sou{koleoptero}. (\sou{zool.)}\cel{3}Insekto kun quar ali, de qui la du supera}
\l{\cel{3}(la elitri), harda, dika, esas quaza etuyi por la du altri.}
\l{\cel{3}- DEFIRSL.}
\l{}
\l{\sou{kolero}. (\sou{patol}.) Morbo epidemiala, quan karakterizas grava pertur-}
\l{\cel{3}beso intestinala, kun loldesko interna e kontrakto spasmatra}
\l{\cel{3}di la membri. - DEFIRS.}
\l{}
\l{\sou{kolerino.}\cel{2}(\sou{patol.})\cel{2}Diareo fortega qua esas (ulakaze) simptomo}
\l{\cel{3}qua anuncas kolero. - DEFS.}
\l{}
\l{\sou{koliar}. (\sou{trans}.) Deprenar de la stipo (floro, frukto, ramo).}
\l{\cel{3}- FIS.}
\l{}
\l{\sou{koliaro}. I. Ringo qua cirkondas la kolo : l. Ornivo por la kolo.}
\l{\cel{3}2. Ringo ek metalo, ek ledro, quan onu m etis cirkum la kolo}
\l{\cel{3}di la sklavi, quan onu metas cirkum la kolo di la animali por}
\l{\cel{3}ligar li ad ulo. 3. Kolo-harneso di la tiro-bestii. 4. Ringo}
\l{\cel{3}qua cirkondas ula objekti. - DEFIS.}
\l{}
\l{\sou{kolibrio}. (\sou{zool.}) Ucelo di la regioni tropikala di Amerika, kun}
\l{\cel{3}beko arkatra,kolori briloza, la maxim mikra ek la uceli konoc-}
\l{\cel{3}ata. - L.\cel{1}\sou{trochilus}. - DEFIRS.}
\l{}
\l{\sou{kolikar}. (\sou{netrans.})\cel{2}(\sou{patol.)}\cel{2}Dolorar en la intestini.-DEFIRS.}
\l{}
\l{\sou{kolimatoro}. (\sou{optiko)}. Aparato uzata en la artilrio por vizar lor}
\l{\cel{3}la apunto di kanono. - DEF.}
\l{}
\l{\sou{kolimbo}. (\sou{zool.}) Ucelo aquala, qua povas permanar dum multa tem-}
\l{\cel{3}po sur la aquo-surfaco. - L.\cel{1}\sou{colymbus}. - L.}
\l{}
\l{\sou{kolino}. Ter-altajo (min kam 500-metro) qua dominacas la plana}
\l{\cel{3}lando. - FIS.}
\l{}
\l{\sou{kolirio}.\cel{2}(\sou{medicino)}.\cel{2}Medikamento (aquo, pomado, e c.) destin-}
\l{\cel{3}ata ad aplikesor sur la konjunktivo di la okulo. - EF.}
\l{}
\l{\sou{kolizionar}.\cel{2}(\sou{netrans}.) (\sou{aludante du korpi}). Intershokar. -}
\l{\cel{3}DEFIRS.}
\l{}
\l{\sou{kolkotaro}. (\sou{kemio}).\cel{2}Peroxido reda de fero, quan onu obtenas per}
\l{\cel{3}kalcinar la fer-sulfato, e quan onu uzas por polisar la spe-}
\l{\cel{3}guli. - DEFS.}
\l{}
\l{kolmo. La grado maxim alta an qua arivas ulo.}
}
